% -*- coding: utf-8 -*-
% !TEX program = xelatex
\documentclass[plain,noanswer,medmath]{../../template/jnuexam}
\usepackage{../../template/kysx}

\begin{document}


\examtitle{1998}{2}


\loadproblems{1}{1998P1}%载入试卷一


\makepart{填空题}{本题共5小题，每小题3分，满分15分}


\useproblem{1}{1}{1}%【同试卷一第一(1)题】


\begin{problem}
曲线$y = - x^3 + x^2 + 2x$与$x$轴所围成的图形的面积$A =$ \fillin{}．
\end{problem}


\begin{problem}
$\int_{}^{} \frac{\ln(\sin x)}{\sin^2 x}\dx =$ \fillin{}．
\end{problem}


%【类似试卷一第二(1)题】
\begin{problem}
设$f(x)$连续，则$\frac{\d}{\dx}\int_0^x tf(x^2 - t^2)\dt =$ \fillin{}．
\end{problem}


\begin{problem}
曲线$y = x\ln \left(\e + \frac{1}{x} \right)$（$x > 0$）的渐进线方程为 \fillin{}．
\end{problem}


\makepart{选择题}{本题共5小题，每小题3分，满分15分}


\begin{problem}
设数列$x_n$与$y_n$满足$\lim_{x \to \infty} x_n y_n = 0$，则下列断言正确的是\pickout{}
\begin{abcd}
\item 若$x_n$发散，则$y_n$必发散．
\item 若$x_n$无界，则$y_n$必有界．
\item 若$x_n$有界，则$y_n$必为无穷小．
\item 若$\frac{1}{x_n}$为无穷小，则$y_n$必为无穷小．
\end{abcd}
\end{problem}


\useproblem{1}{2}{2}%【同试卷一第二(2)题】


\useproblem{1}{2}{3}%【同试卷一第二(3)题】


\begin{problem}
设函数$f(x)$在$x = a$的某个邻域内连续，且$f(a)$为其极大值，
则存在$\delta > 0$，当$x \in \left(a - \delta ,a + \delta \right)$时，必有\pickout{}
\begin{abcd}
\item $(x - a)[f(x) - f(a)] \ge 0$．
\item $(x - a)[f(x) - f(a)] \le 0$．
\item $\lim_{t \to a} \frac{f(t) - f(x)}{(t - x)^2} \ge 0\;\left(x \ne a \right)$．
\item $\lim_{t \to a} \frac{f(t) - f(x)}{(t - x)^2} \le 0\;\left(x \ne a \right)$．
\end{abcd}
\end{problem}


\begin{problem}
设$A$是任一$n$（$n \ge 3$）阶方程，$A^*$是其伴随矩阵，又$k$为常数，
且$k \ne 0, \pm 1$，则必有$(kA)^* =$\pickout{}
\begin{abcd}
\item $kA^*$．
\item $k^{n - 1}A^*$．
\item $k^n A^*$．
\item $k^{-1}A^*$．
\end{abcd}
\end{problem}


\makepart{}{本题满分5分}

\begin{problem}
求函数$f(x) = \left(1 + x \right)^{\frac{x}{\tan(x - \pi /4)}}$．
在区间$\left(0,2\pi \right)$内的间断点，并判断其类型．
\end{problem}


\makepart{}{本题满分5分}

\begin{problem}
确定常数$a,b,c$的值，
使$\lim_{x \to 0}\frac{ax-\sin x}{\int_b^x\frac{\ln(1 + t^3)}{t}\dt} = c$（$c \ne 0$）．
\end{problem}


\makepart{}{本题满分5分}

\begin{problem}
利用代换$y = \frac{u}{\cos x}$将方程$y''\cos x - 2y'\sin x + 3y\cos x = \e^x$化简，并求出原方程的通解．
\end{problem}


\makepart{}{本题满分6分}

\begin{problem}
计算积分$\int_{\frac{1}{2}}^{\frac{3}{2}} \frac{\dx}{\sqrt{\left| x - x^2 \right|}}$．
\end{problem}


\makepart{}{本题满分6分}%【同试卷一第五题】

\useproblem{1}{5}{0}


\makepart{}{本题满分8分}%【同试卷一第九题】

\useproblem{1}{9}{0}


\makepart{}{本题满分8分}

\begin{problem}
设有曲线$y = \sqrt{x - 1}$，过原点作其切线，
求由此曲线、切线及$x$轴围成的平面图形绕$x$轴旋转一周所得到的旋转体的表面积．
\end{problem}


\makepart{}{本题满分8分}

\begin{problem}
设$y = y(x)$是一向上凸的连续曲线，其上任意一点$(x,y)$处的曲率为$\frac{1}{\sqrt{1 + y'^2}}$，
且此曲线上点$(0,1)$处的切线方程为$y = x + 1$，求该曲线的方程，并求函数$y = y(x)$的极值．
\end{problem}


\makepart{}{本题满分6分}

\begin{problem}
设$x \in \left(0,1 \right)$，证明：
\begin{enumerate*}
\item $(1 + x) \ln^2(1 + x) < x^2$；
\item $\frac{1}{\ln 2} - 1 < \frac{1}{\ln(1 + x)} - \frac{1}{x} < \frac{1}{2}$．
\end{enumerate*}
\end{problem}


\makepart{}{本题满分5分}

\begin{problem}
设$\left(2E - C^{-1} B \right) A^T = C^{-1}$，其中$E$是4阶单位矩阵，$A^T$是$4$阶矩阵$A$的转置矩阵，
$B = \begin{pmatrix*}[r]
  1 & 2 & -3 & -2 \\
  0 & 1 &  2 & -3 \\
  0 & 0 &  1 & 2 \\
  0 & 0 &  0 & 1
\end{pmatrix*}$，$C = \begin{pmatrix}
  1 & 2 & 0 & 1 \\
  0 & 1 & 2 & 0 \\
  0 & 0 & 1 & 2 \\
  0 & 0 & 0 & 1
\end{pmatrix}$，求$A$．
\end{problem}


\makepart{}{本题满分8分}

\begin{problem}
已知$\alpha_1 = (1,4,0,2)^T$, $\alpha_2 = (2,7,1,3)^T$, $\alpha_3 = (0,1, - 1,a)^T$, 
$\beta = (3,10,b,4)^T$．问：
\begin{enumerate}
\item $a$, $b$取何值时，$\beta$不能由$\alpha_1$, $\alpha_2$, $\alpha_3$线性表示？
\item $a$, $b$取何值时，$\beta$可由$\alpha_1$, $\alpha_2$, $\alpha_3$线性表示？并写出表达式．
\end{enumerate}
\end{problem}


\end{document}
